\section{Secrets, identity, and authorization}

\subsection{Identity model}

Human access uses enterprise SSO and multi-factor authentication. Service-to-service
access uses mutually authenticated TLS identities. A job obtains a unique OIDC workload
identity whose claims include organization, repository, source revision, run, job,
attempt, environment, plan digest, runner pool, and token audience. Default token
lifetime is five minutes and cannot exceed the remaining job lease.

Cloud and environment access uses federation from this identity; long-lived cloud
keys are prohibited. Token exchange policy matches exact issuer, audience, repository,
branch or tag rule, environment, and protected-pool claim.

\subsection{Secret delivery}

The pipeline plan carries secret references, never values. At step start the runner
requests an access grant from the control plane. After a fresh authorization decision,
the runner exchanges its workload identity directly with the secret manager. Secret
values are delivered in memory or a job-local in-memory filesystem and are removed
after the step. Values are not included in environment dumps, command arguments,
process titles, crash reports, caches, or artifacts.

Secret access from untrusted fork code is denied. Environment secrets are available
only after environment policy succeeds. Rotation does not require rebuilding runner
images. Every read creates an audit event with secret identifier, not value.

\subsection{Authorization model}

Access control combines role-based permissions with attributes such as organization,
repository, environment, source ref, event type, actor, runner pool, and risk score.
Deny rules override allow rules. Policy evaluation returns a decision identifier,
matched policy version, reasons, obligations, and expiry.

\begin{table}[htbp]
\centering
\caption{Minimum authorization roles}
\label{tab:roles}
\begin{tabularx}{\textwidth}{@{}L{0.22\textwidth}Y Y@{}}
  \toprule
  \thead{Role} & \thead{Permitted actions} & \thead{Explicit exclusions} \\
  \midrule
  Viewer & View authorized runs, metadata, and redacted logs & Retry, cancel, approve, administer \\
  Developer & Trigger, cancel own run, retry non-protected jobs & Change policy, approve own deployment \\
  Maintainer & Manage repository pipelines and ordinary schedules & Administer platform or bypass environment policy \\
  Environment approver & Approve or reject eligible deployments & Modify the pending artifact or approve own request if separation is required \\
  Organization admin & Manage repositories, quotas, pools allowed to the organization & Access platform break-glass or rewrite audit \\
  Platform operator & Operate capacity and services & Read secrets or approve application deployments by default \\
  Auditor & Read retained audit, policy decisions, and provenance & Mutate any production resource \\
  \bottomrule
\end{tabularx}
\end{table}

\subsection{Approvals and separation of duties}

Approval rules specify eligible groups, minimum count, self-approval policy, expiry,
and whether a new commit, artifact, plan, policy, or vulnerability result invalidates
approval. The approval record binds the exact environment, plan digest, artifact
digest, actor, decision, and timestamp. Production defaults to two eligible approvers,
no self-approval, and 24-hour expiry. Emergency access is time-bound, reason-coded,
independently alerted, and retrospectively reviewed.

% -----------------------------------------------------------------------------
